\section{Source passages and artifact scope}
\label{app:source-evidence}

All source equation and page locators refer to the PDF of
arXiv:2607.07106v1~\cite{original}. The supplied ten-page PDF is byte-identical
to a fresh versioned arXiv download: 7,847,791 bytes, SHA-256
\begin{quote}\footnotesize\ttfamily
fc5344e2ba8f4db98baf90d805bb6b7c\\
e8835c32ea7bf54bba4d62f405b955a3.
\end{quote}
The official history checked on 28 October 2026 lists v1 submitted
8 July 2026 at 07:43:25 UTC, with no later revision or journal reference.
The PDF displays 9 July; the experimental HTML displays 24 August. The original
\texttt{apssamp.tex}, line 95, uses \verb|\date{\today}|, so separately compiled
display dates are not version identifiers. No source author has been contacted,
and no author confirmation is claimed.

The separately retrieved
\href{https://arxiv.org/src/2607.07106v1}{versioned source archive} has
7,439,982 bytes and SHA-256
\begin{quote}\footnotesize\ttfamily
4f9abd5f41ffca0f1879981429165ff1\\
a989276792c812e860f57cb5038c3b2e.
\end{quote}
The repository's
\href{https://github.com/Teylersf/counterstreaming-heatflux-audit/tree/main/source-excerpts}{selective excerpts}
contain six attributed files and 17 fragments, with original TeX line and
character positions. A standalone verifier checks the archive and source
hashes and reconstructs every excerpt without executing the source or making
network requests. The full PDF and source archive are not redistributed.

\paragraph{The stated initialization.}
Appendix A, p.~6, states that the electron fluid velocities are initialized from
\[
 \epsilon_{ijk}\partial_j B_k=J_i=q_e n_0\overline v_i.
\]
Appendix C, p.~8, states that the fluid simulations inherit the PIC domain,
grid, boundary conditions and magnetic-field initial condition exactly. It
specifies the common initial drift
\[
 \overline v_{1i}=\overline v_{2i}=\overline v_i
\]
and thermal speed $0.111803c$ for the electron fluids. These are the
matching statements tested in the main text. Original TeX locators are
lines 225, 244 and 280; the printed seed is Eq.~(A3), lines 238--240.

The current public \texttt{udist} documentation describes \texttt{ufl} and
\texttt{spatial\_ufl} as fluid momentum in proper-velocity units and gives
relativistic momentum addition as the default~\cite{osirisudist}. This documents
the input-language interpretation; it does not identify the authors' exact
historical executable. For the archived diagnostics, the definition used here
comes directly from the authors' Zenodo deposit description~\cite{dataset}:
\[
 \mathrm{vfl}_i=\frac1n\int f v_i\,\dd^3v,
 \qquad
 (v_i v_j)_{\rm diagnostic}=\int f v_i v_j\,\dd^3v.
\]
The deck and diagnostic checks, rather than an assumed private code revision,
support the initial-moment comparison.

\paragraph{The density-transport explanation.}
The main text, p.~5, attributes the one-electron-fluid density deficit to the
following mechanism:
\begin{quote}\small
``at the mid-plane $\overline v_y=0$ by symmetry, blocking 1-e-Fluid transport
into the sheet''.
\end{quote}
The complete argument is in original TeX line 191. The symmetry-plane
continuity calculation assesses this causal explanation; the plotted density
deficit remains empirical evidence from the reported benchmark.

\paragraph{The implemented closure and instability claim.}
On p.~4 the source specifies the finite local model
\[
 \partial_m q_{ijm}=\kappa v_{th}(p_{ij}-\delta_{ij}p),
 \qquad p=\tr(p_{ij})/3,\qquad \kappa=0.1d_e^{-1},
\]
with an isotropization timescale of about $90\omega_{pe}^{-1}$
(original TeX line 179). Figure~3's caption, p.~4, explicitly omits its
contribution from the displayed fluid heat flux and notes that only its
divergence is constrained (line 166). The finite term is therefore a disclosed
part of the model. The longitudinal calculation in this study tests that
finite model, rather than silently replacing it by its zero-coefficient limit.
On p.~5 the source also asserts that the two-electron-fluid model can capture
linear growth because the instabilities are due to counter-streaming
(line 207). The general importance of closure choice for instability fidelity
is established prior context~\cite{ng2019}.

\paragraph{Accuracy metric and retained qualifications.}
The color in source Fig.~4, p.~5, is the best-time score
\[
 \max_t\left[
 1-\frac{\langle|q_{xxx}-q_{CSxxx}|\rangle_{xy}}
 {\langle|q_{xxx}|\rangle_{xy}}
 \right]
\]
(original TeX line 198). Its proposed thermal/drift scaling is explained in
line 202. The present audit distinguishes this empirical best-time diagnostic
from a uniform error guarantee; it does not claim to have rerun the parameter
scan. The source acknowledges later stream folding (p.~5, line 207), separate
fluid color scales and amplitude differences in the full-tensor comparison
(Fig.~5, p.~7, line 250), and kinetic free-expansion limitations
(Appendix C, p.~8, line 293). The exact mixture identities in source
Eqs.~(3)--(4) and (B1)--(B2) are preserved as valid algebra.
